% !TeX root = ../../Book4.tex
% 纲要标题：### 20.2 抽象 t-结构
% 纲要标签：b4:sec:1902
% 纲要位置：计划纲要.md，第 3748 行。

\section{抽象 t-结构}
\label{b4:sec:1902}

截断的核心数据是两类对象之间的正交性和每个对象的分解三角。将这些条件抽象出来，就能在未指定复形代表的三角范畴中恢复上同调函子。


% 纲要内容（仅作写作计划，尚非教材正文）：
%
% * 两个子范畴及正交、分解公理
% * 有界性与非退化性
% * 由截断恢复上同调函子
%

\subsection{t-结构的正交与分解公理}
\label{b4:subsec:1902:01}

标准截断的存在不必靠复形的每一项描述。上节的两类对象、它们之间的态射零性，以及把任意对象分为两部分的三角，已经足以决定截断函子。

\begin{definition}\label{b4:def:t-structure}
设 \(\mathcal D\) 为三角范畴。取包含零对象且在同构下封闭的两个全子范畴 \(\mathcal D^{\le0},\mathcal D^{\ge0}\)，并记 \(\mathcal D^{\le n}=\mathcal D^{\le0}[-n]\)、\(\mathcal D^{\ge n}=\mathcal D^{\ge0}[-n]\)。若满足
\begin{enumerate}
\item \(\mathcal D^{\le0}[1]\subseteq\mathcal D^{\le0}\)，\(\mathcal D^{\ge0}[-1]\subseteq\mathcal D^{\ge0}\)；
\item \(\operatorname{Hom}(U,V)=0\)，其中 \(U\in\mathcal D^{\le0}\)、\(V\in\mathcal D^{\ge1}\)；
\item 每个 \(X\in\mathcal D\) 都能放入好三角 \(U\to X\to V\to U[1]\)，其中 \(U\in\mathcal D^{\le0}\)、\(V\in\mathcal D^{\ge1}\)，
\end{enumerate}
则称这对子范畴为 \(\mathcal D\) 上的一个\term[tjiegou]{t-结构}[t-structure]。记 \(\mathcal D^{[a,b]}=\mathcal D^{\ge a}\cap\mathcal D^{\le b}\)，其中 \(a,b\in\mathbb Z\)。交 \(\mathcal H=\mathcal D^{\le0}\cap\mathcal D^{\ge0}\) 称为它的\term[xin]{心}[heart]。
\end{definition}
\symidx{\(\mathcal D^{\le n},\mathcal D^{\ge n}\)：给定 t-结构的次数子范畴}
\symidx{\(\mathcal D^{[a,b]}\)：给定 t-结构中次数位于 \([a,b]\) 的子范畴}
\symidx{\(\mathcal H\)：给定 t-结构的心}

移位方向由 \(C[1]^i=C^{i+1}\) 决定。若 \(M\) 是交换范畴中的对象，把它看成标准心中的零次复形，则
\[
 H^{-1}(M[1])=M,\qquad H^0(M[0])=M,\qquad H^1(M[-1])=M,
\]
且各对象的其他上同调均为零。因此 \([1]\) 将这一层移到负一次，\([-1]\) 将它移到一次；定义中的 \([-n]\) 正是把零次条件移到第 \(n\) 次。

\ProseParagraphBreak
例如，设 \(k\) 为域，\(\mathcal D=D^b(\mathbf{Vect}_k^{\mathrm{fd}})\)。标准 t-结构的心由集中在零次的有限维向量空间组成。把两个标准子范畴同时平移为 \((\mathcal D^{\le1},\mathcal D^{\ge1})\)，仍得到 t-结构，其心由 \(V[-1]\) 组成。于是同一个 \(k[-1]\) 在标准心之外，却在平移后的心中。这说明“哪些对象作为零次对象”取决于 t-结构的选择，而三角范畴本身并未改变。

\begin{theorem}\label{b4:thm:t-truncation-adjunction}
设 \(\mathcal D\) 为带 t-结构 \((\mathcal D^{\le0},\mathcal D^{\ge0})\) 的三角范畴，\(n\in\mathbb Z\)，并以 \([-n]\) 平移定义 \(\mathcal D^{\le n},\mathcal D^{\ge n}\)。包含 \(\mathcal D^{\le n}\to\mathcal D\) 有右伴随 \(\tau^{\le n}\)，包含 \(\mathcal D^{\ge n+1}\to\mathcal D\) 有左伴随 \(\tau^{\ge n+1}\)。每个对象 \(X\in\mathcal D\) 有截断三角
\[
 \tau^{\le n}X\longrightarrow X\longrightarrow\tau^{\ge n+1}X
 \longrightarrow(\tau^{\le n}X)[1],
\]
它在保持中间对象及两端次数条件的意义下唯一到唯一同构，并对 \(X\) 的态射函子性地变化。
\end{theorem}
\begin{proof}
低端截断要使每个来自 \(T\in\mathcal D^{\le n}\) 的态射 \(T\to X\) 唯一经过 \(\tau^{\le n}X\to X\)；高端截断要使每个指向 \(W\in\mathcal D^{\ge n+1}\) 的态射 \(X\to W\) 唯一经过 \(X\to\tau^{\ge n+1}X\)。换言之，所需的自然双射分别是
\[
\begin{aligned}
 \operatorname{Hom}(T,\tau^{\le n}X)&\xrightarrow{\sim}\operatorname{Hom}(T,X),\\
 \operatorname{Hom}(\tau^{\ge n+1}X,W)&\xrightarrow{\sim}\operatorname{Hom}(X,W).
\end{aligned}
\]
前一式把截断放在 Hom 的第二变量，给出包含函子的右伴随；后一式把截断放在第一变量，给出包含函子的左伴随。

平移后只须处理 \(n=0\)。取分解公理中的三角 \(U\to X\to V\)。对 \(T\in\mathcal D^{\le0}\)，由\cref{b4:prop:triangle-hom-exact}得到的 Hom 长正合列在 \(\operatorname{Hom}(T,U)\to\operatorname{Hom}(T,X)\) 两旁的项为 \(\operatorname{Hom}(T,V[-1])\) 与 \(\operatorname{Hom}(T,V)\)。它们为零，因为 \(V[-1]\in\mathcal D^{\ge2}\subseteq\mathcal D^{\ge1}\)。所以该映射是同构，给出右伴随的泛性质。对 \(W\in\mathcal D^{\ge1}\)，在 \(\operatorname{Hom}(V,W)\to\operatorname{Hom}(X,W)\) 两旁出现 \(\operatorname{Hom}(U[1],W)\)、\(\operatorname{Hom}(U,W)\)，也都为零；这给出左伴随。

泛性质唯一确定两端对象及其结构映射。给定 \(X\to X'\)，两端态射分别由伴随唯一确定；三角公理 TR3 将前两项的交换方块补成三角态射，而所补的第三项因左伴随的唯一性，必须等于已经确定的态射。

对固定的中间对象，先用这些泛性质给出的唯一同构识别两端，便可把两个分解三角写成
\[
\begin{gathered}
U\xrightarrow{i}X\xrightarrow{p}V\xrightarrow{w}U[1],\\
U\xrightarrow{i}X\xrightarrow{p}V\xrightarrow{w'}U[1].
\end{gathered}
\]
TR3 首先给出三角态射 \((1_U,1_X,c)\)，其中 \(c:V\to V\) 满足 \(cp=p\)。于是 \((c-1_V)p=0\)。对第一三角应用反变 Hom 正合列，得到
\[
\operatorname{Hom}(U[1],V)\longrightarrow
\operatorname{Hom}(V,V)\longrightarrow\operatorname{Hom}(X,V).
\]
因此存在 \(a:U[1]\to V\)，使 \(c-1_V=aw\)。但 \(U[1]\in\mathcal D^{\le0}\)、\(V\in\mathcal D^{\ge1}\)，正交性给出 \(a=0\)，故 \(c=1_V\)。最后一个方块随之给出 \(w'=w\)，而同一论证也说明第三分量唯一。由此得到保持中间对象的唯一三角同构；态射上的恒等与复合则由两端伴随的唯一性保证。
\end{proof}

\begin{proposition}\label{b4:prop:t-orthogonals-extensions}
设 \(\mathcal D\) 为带 t-结构 \((\mathcal D^{\le0},\mathcal D^{\ge0})\) 的三角范畴，\(n\in\mathbb Z\)。则 \(\mathcal D^{\le n}\) 恰是 \(\mathcal D^{\ge n+1}\) 的左正交类，\(\mathcal D^{\ge n+1}\) 恰是 \(\mathcal D^{\le n}\) 的右正交类。二者均对好三角中的扩张、有限直和及直和因子封闭。
\end{proposition}
\begin{proof}
若 \(\operatorname{Hom}(T,X)=0\) 对所有 \(T\in\mathcal D^{\le n}\) 成立，伴随同构使 \(\operatorname{Hom}(\tau^{\le n}X,\tau^{\le n}X)=0\)，故 \(\tau^{\le n}X=0\)，截断三角给出 \(X\simeq\tau^{\ge n+1}X\)。反向是正交公理。左正交刻画对偶可得。对任意好三角 \(X\to Y\to Z\)，若 \(X,Z\) 都与指定的一类对象正交，Hom 长正合列表明 \(Y\) 也正交。有限直和与直和因子的正交性则由 Hom 的加性得到。
\end{proof}

\subsection{有界性与非退化性}
\label{b4:subsec:1902:02}

\begin{definition}\label{b4:def:bounded-nondegenerate-t}
设 \(\mathcal D\) 为带 t-结构 \((\mathcal D^{\le0},\mathcal D^{\ge0})\) 的三角范畴。若每个对象 \(X\in\mathcal D\) 都属于某个 \(\mathcal D^{\ge a}\cap\mathcal D^{\le b}\)，其中整数 \(a,b\) 可以依赖 \(X\)，则称该 t-结构\term[youjietjiegou]{有界}[bounded]。若
\[
 \bigcap_n\mathcal D^{\le n}=0,\qquad
 \bigcap_n\mathcal D^{\ge n}=0,
\]
则称它\term[feituihuatjiegou]{非退化}[nondegenerate]。这里等于零表示交中只有零对象的同构类。
\end{definition}

有界性蕴含非退化性。例如若 \(X\) 属于所有 \(\mathcal D^{\le n}\)，又因有界性属于某个 \(\mathcal D^{\ge a}\)，则 \(X\in\mathcal D^{\le a-1}\cap\mathcal D^{\ge a}\)，正交性迫使 \(1_X=0\)。另一侧相同。标准 t-结构在 \(D^b(\mathcal A)\) 上有界，在一般 \(D(\mathcal A)\) 上未必有界，却仍非退化：所有上同调为零的复形在导出范畴中就是零对象。

\ProseParagraphBreak
非退化性是一项实质条件。对任意非零三角范畴 \(\mathcal D\)，取 \((\mathcal D,0)\) 也满足 t-结构公理，分解三角为 \(X\xrightarrow{1}X\to0\)。它的心为零，不能从心看见原范畴的非零对象。

\subsection{由截断恢复上同调函子}
\label{b4:subsec:1902:03}

\begin{proposition}\label{b4:prop:t-truncation-commute}
设 \(\mathcal D\) 为带 t-结构 \((\mathcal D^{\le0},\mathcal D^{\ge0})\) 的三角范畴，\(a,b\in\mathbb Z\) 且 \(a\le b\)。两种截断函子的复合 \(\tau^{\ge a}\tau^{\le b}\) 与 \(\tau^{\le b}\tau^{\ge a}\) 自然同构，所得对象属于 \(\mathcal D^{\ge a}\cap\mathcal D^{\le b}\)。
\end{proposition}
\begin{proof}
要提取 \(X\) 中次数位于 \([a,b]\) 的部分，可以先取 \(B=\tau^{\le b}X\)，再去掉其中低于 \(a\) 的部分。令 \(A=\tau^{\le a-1}X\)。因为 \(A\in\mathcal D^{\le b}\)，右伴随泛性质唯一给出 \(A\to B\to X\)，其复合为 \(A\) 的截断结构态射。对每个 \(T\in\mathcal D^{\le a-1}\)，它诱导的映射 \(\operatorname{Hom}(T,A)\to\operatorname{Hom}(T,B)\) 经两端伴随同构，成为 \(\operatorname{Hom}(T,X)\) 的恒等映射。因此 \(A\to B\) 也满足 \(B\) 在 \(a-1\) 处的低端截断泛性质。

把 \(A\to B\) 补成好三角 \(A\to B\to Y\to A[1]\)。由上述泛性质，把前两项与 \(B\) 的截断三角识别；(TR3) 补出第三项的态射，\cref{b4:lem:pretriangle-hom}的二取三同构判据给出 \(Y\simeq\tau^{\ge a}B\)。所以 \(Y\in\mathcal D^{\ge a}\)。同时 \(B,A[1]\in\mathcal D^{\le b}\)，对旋转三角 \(B\to Y\to A[1]\) 使用\cref{b4:prop:t-orthogonals-extensions}，得到 \(Y\in\mathcal D^{\le b}\)。这说明 \(Y\) 确实只保留所需次数区间。

对复合 \(A\to B\to X\)，两条已有的截断三角分别以 \(\tau^{\ge a}X\)、\(\tau^{\ge b+1}X\) 为第三项。八面体公理将它们与 \(A\to B\) 的三角联系为
\[
 Y\longrightarrow\tau^{\ge a}X\longrightarrow\tau^{\ge b+1}X\longrightarrow Y[1].
\]
由于 \(Y\in\mathcal D^{\le b}\)，最后一个高端对象属于 \(\mathcal D^{\ge b+1}\)，这又是 \(\tau^{\ge a}X\) 在 \(b\) 处的截断三角，故 \(Y\simeq\tau^{\le b}\tau^{\ge a}X\)。两次识别得到所述同构。八面体中的 \(Y\to\tau^{\ge a}X\) 与 \(B\to X\to\tau^{\ge a}X\) 相容，因而也由 \(B\) 的高端截断泛性质唯一确定；各结构态射都对 \(X\) 自然，故所得同构自然。
\end{proof}

两侧截断 \(\tau^{\ge n}\tau^{\le n}X\) 将对象限制在第 \(n\) 层；它属于 \(\mathcal D^{\ge n}\cap\mathcal D^{\le n}\)，还须平移 \([n]\) 才把这一层移回零次的心。因此在未指定复形各项的情况下定义
\[
 H_t^n(X)\coloneqq\bigl(\tau^{\ge n}\tau^{\le n}X\bigr)[n]
                 =H_t^0(X[n])\in\mathcal H.
\]
这称为关于给定 t-结构的\term[tshangtongdiao]{上同调对象}[cohomology object]。例如在标准结构中，\(M[-2]\) 只有二次上同调，取两侧截断后仍是 \(M[-2]\)；再施加 \([2]\)，得到 \(M[0]\)，所以 \(H_t^2(M[-2])=M\) 是心中的对象。下标 \(t\) 用来区别随后出现的新心；标准情形可省去。上述定义先给出取值于加性心的函子；心的交换性和三角诱导的长正合列将在下一节证明。
\symidx{\(H_t^n\)：给定 t-结构的上同调函子}

\ProseParagraphBreak
相邻截断之间有好三角 \(\tau^{\le n-1}X\to\tau^{\le n}X\to H_t^n(X)[-n]\to\)，这是前一证明取 \(a=b=n\) 的情形。有界对象因此能由有限多个上同调对象逐次扩张得到。每一层的连接映射仍是重建对象所需的数据。抽象截断与交换关系的原始表述见\cite[Propositions 1.3.3, 1.3.5]{BBD1982FaisceauxPervers}。
